(1) 轴对称纯 \(\ell=1\) 解可写
\[
\Phi_1=-AE_0r\cos\theta,\quad
\Phi_2=-E_0(Br+C/r^2)\cos\theta,\quad
\Phi_3=-E_0(r-D/r^2)\cos\theta .
\]
在 \(r=a,b\) 均有电势连续 \([\Phi]=0\) 与法向位移连续 \([\varepsilon\partial_r\Phi]=0\)，因为无自由面电荷。外部偶极势系数 \(E_0D\)，故 \(p=4\pi\varepsilon_0\varepsilon_3E_0D\)。

(2) 在核壳界面 \(A=B+C/a^3\)、\(\varepsilon_1A=\varepsilon_2(B-2C/a^3)\)，得到 \(C/a^3=(\varepsilon_2-\varepsilon_1)B/(\varepsilon_1+2\varepsilon_2)\)。再代入外界面方程，定义
\[
\varepsilon_{\rm eff}=\varepsilon_2
\frac{\varepsilon_1+2\varepsilon_2+2f(\varepsilon_1-\varepsilon_2)}
{\varepsilon_1+2\varepsilon_2-f(\varepsilon_1-\varepsilon_2)}.
\]
消元得 \(D=b^3(\varepsilon_{\rm eff}-\varepsilon_3)/(\varepsilon_{\rm eff}+2\varepsilon_3)\)，故
\[
\boxed{\alpha_e=\frac p{E_0}=4\pi\varepsilon_0\varepsilon_3b^3
\frac{\varepsilon_{\rm eff}-\varepsilon_3}{\varepsilon_{\rm eff}+2\varepsilon_3}}.
\]
当 \(f\to0\) 或 \(\varepsilon_1=\varepsilon_2\) 时均退化为均匀介质球的极化率。

(3) 中性包覆要求 \(\alpha_e=0\)，即 \(\varepsilon_{\rm eff}=\varepsilon_3\)，等价于
\[
(\varepsilon_2-\varepsilon_3)(\varepsilon_1+2\varepsilon_2)
+f(\varepsilon_1-\varepsilon_2)(\varepsilon_3+2\varepsilon_2)=0.
\]
均匀远场仅激发 \(\ell=1\) 通道，因此 \(D=0\) 使 \(r>b\) 整个外部静电场严格等于原均匀场，并非仅在无限远处近似相等；内部电场和极化电荷仍可非零。

(4) 由内界面条件 \(A=3\varepsilon_2B/(\varepsilon_1+2\varepsilon_2)\)，其中
\[
B=\frac{3\varepsilon_3}{\varepsilon_{\rm eff}+2\varepsilon_3}
\frac{\varepsilon_1+2\varepsilon_2}{\varepsilon_1+2\varepsilon_2-f(\varepsilon_1-\varepsilon_2)}.
\]
核内 \(\boldsymbol E_1=AE_0\boldsymbol e_z\)；壳层 \(E_{2r}=E_0(B-2C/r^3)\cos\theta\)、\(E_{2\theta}=-E_0(B+C/r^3)\sin\theta\)。线性响应且外部电源固定 \(E_0\) 时，插入球体的极化自由能 \(\Delta F=-\int_0^{E_0}p(E)dE=-\alpha_eE_0^2/2\)，包含电源的做功；固定自由电荷边界条件下必须变换热力学势，不能照搬此负号。
